\section{Uncertainty quantification: three routes}
\label{sec:inference}

A point estimate without standard errors is half an answer. The library
implements three inference routes with deliberately different assumptions
and costs, summarized in Table~\ref{tab:inference-routes}. Offering all
three is itself a design decision: comparing their finite-sample coverage
across jump families is one of the research questions this library was
built to answer (Section~\ref{sec:closing}).

\begin{table}[ht]
\centering
\small
\begin{tabular}{@{}p{5.6cm}p{4.5cm}l@{}}
\toprule
Route (method) & Assumes & Cost \\
\midrule
Wald\newline {\footnotesize\code{estimate\_wald\_standard\_errors}} &
  asymptotic normality, regular MLE & one Hessian \\
\addlinespace
Profile\newline {\footnotesize\code{estimate\_standard\_errors}} &
  Wilks' $\chi^2$ asymptotics & re-optimizations \\
\addlinespace
Bootstrap\newline {\footnotesize\code{estimate\_bootstrap\_standard\_errors}} &
  correct model specification & $B$ full re-fits \\
\bottomrule
\end{tabular}
\caption{The three inference routes. Cost grows downward; robustness to
boundary and non-quadratic likelihoods grows the same way.}
\label{tab:inference-routes}
\end{table}

\subsection{Wald intervals from the observed information}
\label{sec:wald}

Classical first-order asymptotics: under regularity conditions,
$\hat\thetavec \approx \mathcal N\bigl(\thetavec_0,\ I(\hat\thetavec)^{-1}\bigr)$
where
\begin{equation}
  I(\hat\thetavec) \;=\; -\nabla^2 \loglik(\hat\thetavec)
  \label{eq:observed-info}
\end{equation}
is the \emph{observed} Fisher information --- the Hessian of the negative
log-likelihood at the estimate. Standard errors are the square roots of
the diagonal of $I^{-1}$, and intervals are
$\hat\theta_r \pm z_{1-\alpha/2}\,\widehat{\mathrm{SE}}_r$.

\begin{decision}[Observed, not expected, information]
Following \citet{EfronHinkley1978}, the library uses the observed
information rather than the expected Fisher information: for this mixture
likelihood the expectation has no closed form and would itself require
numerical integration, and Efron and Hinkley argue the observed
information is the more relevant conditioning quantity anyway.
\end{decision}

Since the code's objective is already the negative log-likelihood, its
Hessian is exactly \eqref{eq:observed-info}, built by central finite
differences: on the diagonal
\begin{equation}
  \frac{\partial^2 f}{\partial \theta_r^2}
  \;\approx\;
  \frac{f(\hat\thetavec + h_r e_r) - 2 f(\hat\thetavec)
        + f(\hat\thetavec - h_r e_r)}{h_r^2},
  \label{eq:fd-diag}
\end{equation}
and off the diagonal the four-point mixed formula
\begin{equation}
  \frac{\partial^2 f}{\partial \theta_r \partial \theta_s}
  \;\approx\;
  \frac{f_{++} - f_{+-} - f_{-+} + f_{--}}{4 h_r h_s},
  \label{eq:fd-mixed}
\end{equation}
where $f_{\pm\pm} = f(\hat\thetavec \pm h_r e_r \pm h_s e_s)$.

\begin{decision}[Bound-aware step sizes]
\label{dec:fd-steps}
The step for parameter $r$ is $h_r = 10^{-4} \max(|\hat\theta_r|, 10^{-2})$
--- relative, with an absolute floor so parameters estimated near zero
(drift, jump location) still get a usable step. Crucially, each step is
then shrunk to at most a quarter of the distance to the nearest bound:
the objective returns $+\infty$ outside the bounds
(Decision~\ref{dec:penalize}), and a single infinite evaluation would
destroy the difference quotients \eqref{eq:fd-diag}--\eqref{eq:fd-mixed}.
Entries that still cannot be evaluated (a parameter pinned exactly to a
bound) are recorded as \code{NaN}, and the inversion is guarded: if the
Hessian is non-finite or singular, every Wald SE is reported as
\code{NaN} rather than silently wrong. This is not an edge case ---
$\hat p$ near $0$ or a jump scale collapsing toward its lower bound are
exactly the interesting regimes.
\end{decision}

\subsection{Profile likelihood intervals}
\label{sec:profile}

The Wald construction forces symmetric intervals and trusts a quadratic
approximation of the log-likelihood. The profile likelihood does neither.
For a scalar parameter of interest $\theta_r$, define
\begin{equation}
  \loglik_p(v)
  \;=\;
  \max_{\thetavec\,:\,\theta_r = v} \loglik(\thetavec),
  \label{eq:profile}
\end{equation}
the log-likelihood maximized over all \emph{other} parameters with
$\theta_r$ pinned at $v$. By Wilks' theorem \citep{Wilks1938}, under the
true value $2[\loglik(\hat\thetavec) - \loglik_p(\theta_{r,0})]
\rightsquigarrow \chi^2_1$, so
\begin{equation}
  \mathrm{CI}_{1-\alpha}
  \;=\;
  \Bigl\{ v \;:\; \loglik_p(v) \;\ge\;
  \underbrace{\loglik(\hat\thetavec) - \tfrac12 \chi^2_{1,1-\alpha}}_{\text{threshold}}
  \Bigr\}
  \label{eq:profile-ci}
\end{equation}
is an asymptotic $(1-\alpha)$ confidence set that follows the likelihood's
actual shape --- asymmetric where the likelihood is asymmetric, respecting
parameter bounds by construction.

Each profile point \eqref{eq:profile} is a constrained re-optimization
(L-BFGS-B over the remaining parameters, warm-started at the MLE), so the
practical question is \emph{where} to spend those evaluations. The
library's algorithm, per parameter:

\begin{enumerate}
  \item \textbf{Adaptive outward search.} Starting from the MLE with step
        $\max(0.01\,|\hat\theta_r|,\ 10^{-4})$, walk right, doubling the
        step each time, evaluating $\loglik_p$ at each point, until the
        profile drops below the threshold in \eqref{eq:profile-ci}, a
        parameter bound is hit (evaluate just inside it), or 15 steps are
        exhausted. Repeat leftward. Geometric doubling reaches a crossing
        $k$ standard errors away in $O(\log k)$ evaluations --- important
        because parameters like $p$ can have profiles an order of
        magnitude wider than their Wald SE suggests.
  \item \textbf{Infill.} If fewer than \code{n\_points} evaluations
        accumulated, bisect the largest remaining gap between consecutive
        evaluated points, repeatedly --- concentrating resolution where
        the profile is least known.
  \item \textbf{Interval endpoints by interpolation.} The CI endpoints are
        where $\loglik_p$ crosses the threshold: found by linear
        interpolation between the bracketing evaluations on each side. If
        the profile never drops below the threshold before a bound (e.g.\
        $p$ against $0$), the interval is truncated at the searched range
        --- reproducing the one-sided behavior a boundary parameter should
        have.
  \item \textbf{Standard error from local curvature.} Near the MLE,
        $\loglik_p(v) \approx \loglik(\hat\thetavec) -
        (v - \hat\theta_r)^2 / (2\,\mathrm{SE}^2)$: a parabola whose
        leading coefficient $a$ satisfies
        $\mathrm{SE} = \sqrt{-1/(2a)}$. The code fits that parabola by
        least squares through the (up to five) highest-likelihood profile
        points. If the fit is not concave ($a \ge 0$, a genuinely
        non-quadratic profile), the fallback derives the SE from the
        interval width instead, $\mathrm{SE} \approx
        (\mathrm{CI}_{\mathrm{high}} - \mathrm{CI}_{\mathrm{low}}) /
        (2 z_{1-\alpha/2})$ --- and only when the endpoints genuinely
        crossed the threshold rather than being truncated at the searched
        range.
\end{enumerate}

\begin{lstlisting}[caption={The profiled objective: pin $\theta_r$, optimize the rest (excerpt).}]
def optimize_profile(val):
    def neg_log_lik_profile(free_params):
        full_params = np.zeros(len(self._param_names))
        full_params[idx] = val          # parameter of interest, pinned
        free_idx = 0
        for j in range(len(self._param_names)):
            if j != idx:
                full_params[j] = free_params[free_idx]
                free_idx += 1
        return self.log_likelihood(full_params)

    res = minimize(fun=neg_log_lik_profile, x0=free_init,
                   method="L-BFGS-B", bounds=free_bounds,
                   options={"maxiter": 100, "ftol": 1e-6})
    return -res.fun if res.success and np.isfinite(res.fun) else -np.inf
\end{lstlisting}

The companion method \code{plot\_profiles()} draws each profile with its
threshold line --- the visual version of \eqref{eq:profile-ci}, and in
practice the fastest way to \emph{see} identifiability problems: a flat
profile is a parameter the data cannot pin down.

\subsection{Parametric bootstrap}
\label{sec:bootstrap}

The third route drops the asymptotics altogether and pays with
computation. Treat the fitted model as the truth; simulate $B$ datasets of
the same length $n$ and spacing $\Delta t$ from it; re-estimate
$\thetavec$ on each; and read the sampling distribution of the estimator
directly off the $B$ replicate estimates
$\hat\thetavec^{(1)}, \dots, \hat\thetavec^{(B)}$
\citep{DavisonHinkley1997}:
\begin{equation}
  \widehat{\mathrm{SE}}_r = \operatorname{sd}\bigl(\hat\theta_r^{(b)}\bigr),
  \qquad
  \mathrm{CI}_{1-\alpha} =
  \Bigl[ \hat\theta_r^{(\lceil B\alpha/2 \rceil)},\
         \hat\theta_r^{(\lceil B(1-\alpha/2) \rceil)} \Bigr]
  \quad \text{(percentile interval)}.
  \label{eq:bootstrap}
\end{equation}
No normality, no quadratic approximation, and boundary effects are
represented honestly: if $\hat p$ piles up near zero across replicates,
the percentile interval shows it.

Implementation decisions worth recording: replicate $b$ is simulated with
seed $\texttt{seed} + b$, making the entire bootstrap reproducible from
one integer while keeping replicates independent; re-fits that fail to
converge are discarded (their count is reported as
\code{n\_successful}), since averaging over failed optimizations would
contaminate \eqref{eq:bootstrap} with optimizer artifacts; the SE uses the
$\mathrm{ddof}=1$ sample standard deviation; and the raw replicate matrix
is stored in the results dictionary so users can compute BCa or basic
intervals, correlations between parameters, or replicate diagnostics
without re-running the loop.

\begin{lstlisting}[caption={The bootstrap loop (excerpt).}]
model = JumpDiffusionModel(jump_distribution=jump_dist, **fitted_params)
for b in range(n_bootstrap):
    rep_seed = None if seed is None else seed + b
    _, path, _ = model.simulate(T=horizon, n_steps=n_steps, x0=0.0,
                                seed=rep_seed)
    increments = np.diff(path)
    estimator = JumpDiffusionEstimator(increments, self.dt,
                                       jump_distribution=jump_dist)
    rep_result = estimator.estimate(**estimate_kwargs)
    if rep_result["convergence"]:
        replicates.append([rep_result["parameters"][name]
                           for name in self._param_names])
\end{lstlisting}

All three routes converge into one view: \code{summary()} returns a tidy
table with one row per parameter and one column block per computed route,
and \code{diagnostics()} prints it with convergence and data statistics
--- making the disagreement between routes (the interesting part) visible
at a glance.

\code{diagnostics()} also compares the sample mean and variance of the
increments against the ones the fitted model \emph{implies}. Under the
Bernoulli mixture \eqref{eq:bernoulli-increment} these are exact:
\begin{equation}
  \E[\Delta X] = \mu\,\Delta t + p\,\E[J],
  \qquad
  \operatorname{Var}[\Delta X]
  = \sigma^2 \Delta t + p\,\E[J^2] - \bigl(p\,\E[J]\bigr)^2,
  \label{eq:incr-moments}
\end{equation}
with $\E[J]$ and $\operatorname{Var}[J]$ supplied by the jump law's
generic numerical moments (Section~\ref{sec:contract}).

\begin{remark}[The jump term in the mean is not optional]
An earlier version of \code{diagnostics()} printed $\mu\,\Delta t$ alone
as the theoretical mean. With symmetric jumps that omission is invisible
($\E[J]=0$), but the 2026 audit measured its effect on a skew-normal fit:
a reported theoretical mean of $-0.000019$ against an empirical mean of
$0.011798$ --- suggesting a gross model--data mismatch --- when the
model's actual implied mean, $\mu\Delta t + p\,\E[J] = 0.011778$, agreed
with the data almost exactly. For \emph{asymmetric} jump laws, the whole
point of this library, the $p\,\E[J]$ term routinely dominates
$\mu\,\Delta t$; equation \eqref{eq:incr-moments} is now computed in
full.
\end{remark}
